% §7 Threats to validity (lead). Sources: docs/11 §6 (long form) and docs/04 §11.
\section{Threats to Validity}
\label{sec:threats}

\paragraph{Internal validity}
Harness differences are large in this field: BARO's AC@1 on RE2-TT is 0.67 in RCAEval's harness and 0.32 in PRISM's~\cite{rcaeval2024,prism2026}.
We therefore run every baseline in one harness with a pinned RCAEval commit, the same inputs and the same denominators, compare only against our own re-runs, and report reproduction gaps; our numbers for published methods may still differ from the published ones.
The alert-count baseline \code{simplerca} is a re-implementation from its paper and may differ from the original.
We wrote both the predicates and the verifier, so the trials encode our assumptions; for example, T1 expects the root to show an internal anomaly, which network faults often lack.
All parameters are set a priori in configuration with a cited source and frozen before testing, and we report per-fault-class results, including any high abstention on network faults, rather than tuning it away.
Some defaults came from papers evaluated on RCAEval itself, so some optimism may remain; we treat RE1 as development data and report RE2 and RE3 as test data.
Repetitions of a fault on the same service are not independent, so splits are made by cluster.
The lead author built the system and runs the evaluation; the pre-registration tag, the list of development cases and the deviations log bound that risk but do not remove it.

\paragraph{External validity}
RCAEval's three systems are small demo applications with single injected faults, and real incidents differ.
Even the retrieval of statistical methods, which is 98--100\% on microservice benchmarks, falls to 35--64\% on industrial control data~\cite{decomprca2026}.
Nezha and VeylLab, if their adapters are ready in time, add sources but no production data, and laptop-collected data carries host contention, a desktop virtual machine and a single-node cluster, which we record and report.
The change logs of RCAEval-C are synthetic, and its generated true change encodes the injected service by construction, so RCAEval-C tests how gullible the change channel is, not whether it helps on real change logs, which are noisier; real rollout and flag changes in VeylLab would be the backup evidence.
Conformal coverage assumes exchangeability, which fails across systems; leave-one-system-out coverage is reported as an empirical result.

\paragraph{Construct validity}
Service-level AC@1 ignores the propagation path and the fault type, and LLM agents identify faulty services better than faulty edges~\cite{openrca2026v2}; we also report fault-type and indicator accuracy, but RCAEval has no path ground truth.
Proposer lift depends on the predicate library: the LLM can only help through predicates we wrote, so a zero lift is a finding about this library and model, not about LLMs in general, and we report which predicates the model requested.
RCAEval normal windows are short pre-injection slices that may contain the tail of a previous fault (\S\ref{sec:splits}), so we also report \NOFAULT on VeylLab's no-fault captures if they exist.
RCAEval has known label issues, such as a case whose folder names the wrong code fault~\cite{rcaevalissue43}, missing logs or traces in some cases, and repeated injections; we report per-fault-class results and list the known issues.

\paragraph{Conclusion validity}
Power is low: according to a pre-registered preprint, about 7 points of AC@1 can be detected on 270 cases and about 14 on 90~\cite{hussain2026prereg}, and clustering raises both, so many per-system comparisons will be inconclusive; we use clustered paired tests with Holm correction and the words ``not detected''.
Split conformal needs at least $1/\alpha - 1 = 9$ calibration cases per stratum, so we calibrate per benchmark family, not per family and telemetry modality.

\paragraph{LLM-specific threats}
LLM inference is not bit-exactly reproducible across machines~\cite{thinkingmachinesnondeterminism}; Tier-1 replay from recorded cassettes guarantees the verdict, Tier-2 regeneration is measured, and three seeds expose seed variance.
RCAEval and Nezha have been public since 2023--2024 and may be in pretraining data; no study we found measures this directly.
Anonymising service names cost one LLM pipeline 0.20--0.21 AC@1 on RE3 in a pre-registered preprint~\cite{hussain2026prereg}; our probe (H4d) is optional, exploratory and has no pre-registered experiment, and the verdict is decided by code in any case.
Scores published for the model are vendor-reported, so we report only our own measurements and pin the weights, quantisation and runtime.

\paragraph{Hardware and process}
CPU speed differs across the team's laptops, so cost numbers are hardware-specific; verdicts do not depend on speed, and we report timings per laptop with its specification.
Three people, each owning one lane of the code, wrote the system with the same AI coding agent, Claude Code.
The lead's agent sessions drafted most pull requests, including draft pull requests in the other two lanes, which the lane owners then reviewed.
Not every change had a second human reviewer: in the weeks before the results freeze, the lead merged several pull requests after their own review, with a posted note asking the lane owner to review after the merge.
Agent-written code may contain subtle bugs, and one agent drafting most of the system makes review less independent; the checks that every change passed are ownership lanes, versioned contracts, golden fixtures, tests written first and a continuous-integration determinism gate.
We report this as a threat to validity, not a result, and make no claim about agent productivity.
